% Generated draft: requires mathematical and source review.
\documentclass[11pt]{article}
\usepackage[margin=25mm]{geometry}
\usepackage{fontspec}
\setmainfont{Latin Modern Roman}
\usepackage{amsmath,amssymb,mathtools,unicode-math}
\setmathfont{Latin Modern Math}
\usepackage{enumitem,needspace,xurl,hyperref}
\setlength{\parindent}{0pt}
\setlength{\parskip}{5pt}
\setlength{\emergencystretch}{3em}
\newcommand{\N}{\mathbb{N}}
\newcommand{\Z}{\mathbb{Z}}
\newcommand{\Q}{\mathbb{Q}}
\newcommand{\R}{\mathbb{R}}
\newcommand{\C}{\mathbb{C}}
\newcommand{\F}{\mathbb{F}}
\newcommand{\A}{\mathbb{A}}
\newcommand{\PP}{\mathbb{P}}
\newcommand{\E}{\mathbb{E}}
\newcommand{\eps}{\varepsilon}
\newcommand{\dd}{\,\mathrm{d}}
\begin{document}
\section*{UnsolvedMath ID: 2308013}
\textbf{Source problem code:} catalogue entry
\subsection{Definitions and mathematical statement}
Under the usual circle-measure interpretation, let \(\mathbb T=\R/(2\pi\Z)\), with Lebesgue measure \(m\). For a finite Borel measure \(\mu\), its Fourier--Stieltjes coefficients are
\[
\widehat\mu(k)=\int_{\mathbb T}e^{-ik\theta}\,\dd\mu(\theta),\qquad k\in\Z.
\]
For a positive measure, singularity means that a Borel set \(E\) satisfies \(m(E)=0\) and \(\mu(\mathbb T\setminus E)=0\). Holland's question, in this conventional setting, is whether there exists a nonzero finite singular measure \(\mu\) belonging to Zygmund's class \(A^*\) such that
\[
\widehat\mu(k)\in[0,\infty)\quad\text{for every }k\in\Z.
\]
Nonnegativity requires real coefficients, not merely nonnegative real parts. Nonzero is needed to avoid the zero-measure triviality; for a nonzero positive measure normalization to total mass one does not change the properties.

The record does not define \(A^*\), the underlying space, or whether measures must be positive. For orientation, a common Zygmund-type convention controls second differences of a distribution primitive \(F\) by \(|F(x+h)-2F(x)+F(x-h)|\le C|h|\); an associated measure convention compares equal adjacent intervals. This is not asserted to be the exact \(A^*\) definition intended by the source. Until that definition and the measure category are supplied, the class-membership condition remains an unresolved essential datum. The supplied OPEN-TRIAGE summary reports open status in the 2018 edition and no definitive resolution located.
\subsection{Short English statement}
Is there a nontrivial singular measure in Zygmund's \(A^*\) class with every Fourier--Stieltjes coefficient real and nonnegative? The class definition, measure category, and underlying space are absent.
\subsection{Sources}
\begin{itemize}
\item Local UnsolvedMath registry, permanent ID \(2308013\), exact source code \texttt{catalogue entry}, ``Research Problems in Function Theory --- Problem 8.13.''
\item Hayman, \emph{Research Problems in Function Theory} (2018), Problem 8.13, attributed to F. Holland; \url{https://arxiv.org/abs/1809.07200}. Source-TeX extraction; recorded access 2026-07-29.
\item Zygmund's class is named without a defining reference or formula. The supplied OPEN-TRIAGE summary provides no separate authoritative current-status citation.
\end{itemize}
\end{document}
